\chapter{Theoretical Background and Literature Survey}

\section{Anatomy of a Rootkit: Kernel Space vs. User Space}

Rootkits are classified by their level of system integration. \textbf{User-space rootkits} operate by hooking shared libraries (e.g., via \texttt{LD\_PRELOAD}) to intercept system calls in user mode. They are easier to deploy but also easier to detect through integrity checking tools \cite{jynx2012}.

\textbf{Kernel-space rootkits} operate at the highest privilege level (Ring 0), directly modifying kernel data structures, system call tables, or loadable kernel modules (LKMs). Once installed, they can hide files, processes, network connections, and even other malware from the operating system itself. Their network communication modules operate through raw sockets or Berkeley Packet Filters, bypassing the entire application networking stack \cite{elastic2022bpfdoor}.

This fundamental distinction is critical: kernel rootkits do not generate application-layer logs, rendering log-based defenses completely ineffective.

% ---- FIGURE: Kernel vs User Space ----
\begin{figure}[h]
\centering
\fbox{\parbox{0.85\textwidth}{\centering\vspace{3cm}\textit{Insert Diagram: Kernel Space vs. User Space architecture showing Ring 0 (kernel rootkits, raw sockets, BPF) vs Ring 3 (user-space rootkits, LD\_PRELOAD hooks, application logs)}\vspace{3cm}}}
\caption{Rootkit Classification: Kernel Space (Ring 0) vs. User Space (Ring 3)}
\label{fig:kernel_vs_user}
\end{figure}

\section{Case Studies of Magic Packet Triggers}

\subsection{BPFDoor}
BPFDoor uses the kernel's Berkeley Packet Filter to passively monitor incoming traffic without binding to any port. It activates upon receiving packets containing magic bytes \texttt{0x7255} (UDP/ICMP) or \texttt{0x5293} (TCP) combined with a hardcoded password (e.g., ``justrobot''). This design allows it to operate on systems with strict firewall rules since BPF captures packets before the firewall processes them \cite{elastic2022bpfdoor, trend2022bpfdoor}.

% ---- FIGURE: BPFDoor Packet Structure ----
\begin{figure}[h]
\centering
\fbox{\parbox{0.85\textwidth}{\centering\vspace{2.5cm}\textit{Insert Diagram: BPFDoor magic packet structure showing IP Header $\rightarrow$ UDP/TCP Header $\rightarrow$ Payload with magic bytes 0x7255 and password field highlighted}\vspace{2.5cm}}}
\caption{BPFDoor Magic Packet Structure (UDP Variant)}
\label{fig:bpfdoor_packet}
\end{figure}

\subsection{Cd00r}
Cd00r sniffs incoming SYN packets using raw sockets and activates a reverse shell when it detects a specific sequence of connection attempts. Modern variants encode activation parameters within TCP options fields, payload prefixes (``Z4vE''), and IP address offsets \cite{fx2000cd00r}.

\subsection{Reptile}
Reptile is notable for its high variability across wild samples. Analysis of four deployments reveals completely different configurations: ports (666, 41302, 2307), magic strings (``smoke'', ``7313F4lc0n4710n''), and passwords that change per deployment, rendering static signatures ineffective \cite{trend2023reptile}.

\subsection{Syslogk and Jynx}
Syslogk activates its Rekoobe payload via specific TCP source ports (59318) and hardcoded keys. Jynx uses simple header-pair matching (\texttt{ACK=0xDEAD}, \texttt{SEQ=0xBEEF}), representing the simplest end of the magic packet complexity spectrum \cite{avast2022syslogk, jynx2012}.

\section{Covert Communication Channels in TCP/IP Headers}

Magic packet rootkits exploit multiple fields within the TCP/IP protocol stack for covert signaling:

\begin{itemize}
    \item \textbf{TCP Options Field:} Rarely inspected by standard IDS, rootkits encode activation values at specific offsets (e.g., cd00r uses offset 0x02 with value 1366).
    \item \textbf{Sequence and ACK Numbers:} Used as covert data channels (Jynx encodes magic values; cd00r encodes attacker IP).
    \item \textbf{Window Size:} Linkpro uses window size 54321 as its activation trigger.
    \item \textbf{Source Port Ranges:} Syslogk uses specific port ranges (63400--63411) for control signaling.
    \item \textbf{Payload Magic Bytes:} BPFDoor and Reptile embed unique byte sequences at the start of the payload.
\end{itemize}

% ---- FIGURE: TCP Header Covert Channels ----
\begin{figure}[h]
\centering
\fbox{\parbox{0.85\textwidth}{\centering\vspace{3cm}\textit{Insert Diagram: Annotated TCP/IP header diagram highlighting the fields exploited by rootkits --- Source Port, Seq Number, Ack Number, Window Size, TCP Options, and Payload fields with arrows pointing to which rootkit uses each}\vspace{3cm}}}
\caption{Covert Channels in TCP/IP Headers Exploited by Rootkits}
\label{fig:tcp_covert_channels}
\end{figure}

These channels exploit the fact that most security tools only inspect destination ports and payload signatures, ignoring the rich feature space available in protocol headers.

\section{The Role of Out-of-Band Monitoring in Modern SOCs}

Out-of-Band (OOB) monitoring analyzes a copy of network traffic (obtained via TAP or SPAN port) rather than inspecting packets inline. This architecture is increasingly adopted in Security Operations Centers because it eliminates the risk of network disruption while providing full packet visibility \cite{paxson1999bro}.

% ---- FIGURE: Inline vs OOB Architecture ----
\begin{figure}[h]
\centering
\fbox{\parbox{0.85\textwidth}{\centering\vspace{3cm}\textit{Insert Diagram: Side-by-side comparison of Inline IDS/NGFW architecture (single point of failure in traffic path) vs. OOB/SPAN architecture (sensor receives copy, production traffic unaffected)}\vspace{3cm}}}
\caption{Inline vs. Out-of-Band IDS Architecture Comparison}
\label{fig:inline_vs_oob}
\end{figure}

The key advantage is decoupling: if the analysis engine crashes or requires extended processing for complex fragmented packets, production traffic is entirely unaffected. The trade-off is that OOB systems cannot drop packets in real time --- they must rely on reactive enforcement mechanisms (firewall API calls, webhook-triggered blocking) to mitigate detected threats.

\section{Review of Machine Learning in Traffic Analysis}

Machine Learning has been applied to network intrusion detection to address the limitations of pure signature approaches \cite{buczak2016mlids}. Decision Trees are particularly suitable for real-time packet triage due to:

\begin{itemize}
    \item \textbf{O(depth) Classification:} Microsecond-scale inference enabling per-packet evaluation.
    \item \textbf{Explainability:} Decision paths can be traced and audited, critical for forensic analysis.
    \item \textbf{Low Overhead:} Minimal CPU and memory footprint compared to deep learning models.
\end{itemize}

Shannon Entropy has also been effectively used for detecting encrypted or compressed payloads. Lyda and Hamrock \cite{lyda2007entropy} demonstrated that entropy analysis can distinguish between normal and obfuscated content with high accuracy, as encrypted data exhibits entropy values approaching the theoretical maximum of 8.0 bits per byte \cite{shannon1948}.

\section{Comparison of Existing Solutions}

\begin{table}[h]
\centering
\caption{Comparative Analysis of Existing Security Solutions}
\label{tab:existing_solutions}
\begin{tabular}{|p{2.5cm}|c|c|c|c|}
\hline
\textbf{Capability} & \textbf{Fail2Ban} & \textbf{Snort} & \textbf{Suricata} & \textbf{eBPF Filters} \\
\hline
Raw Packet Visibility & No & Yes & Yes & Yes \\
\hline
Magic Byte Detection & No & Partial & Partial & Manual \\
\hline
Anti-Evasion (XOR/B64) & No & No & No & No \\
\hline
Fragment Reassembly & No & Yes & Yes & No \\
\hline
ML Pre-Filter & No & No & No & No \\
\hline
Auto Rule Generation & No & No & No & No \\
\hline
Automated Response & Yes & Partial & Partial & No \\
\hline
Zero-Trust Auth & No & No & No & No \\
\hline
\end{tabular}
\end{table}

The comparative analysis reveals that no existing open-source solution provides integrated multi-layer anti-evasion, ML-based triage, automated rule inference, and zero-trust authenticated response specifically optimized for magic packet rootkit detection.
